\subsection{Differentiability}
\D{three equivalent formulations}{
$f'(x_0)=\lim_{h\to0}\frac{f(x_0+h)-f(x_0)}{h}=\lim_{x\to x_0}\frac{f(x)-f(x_0)}{x-x_0}$ (\textbf{two-sided}, finite)
$\iff$ $\exists K\in\Rl:\ f(x)=f(x_0)+K(x-x_0)+o(|x-x_0|)$ as $x\to x_0$, and then $K=f'(x_0)$.\\
\hd{diff.\ \ra\ continuous}; continuous $\not\Rightarrow$ diff.\ ($|x|$ at $0$).
}
\X{Landau $o$ / $O$ --- the classic multiple-choice trap}{
$f=O(g)$: $|f|\le C|g|$ near $x_0$ (\emph{same order or smaller}).\quad
$f=o(g)$: $f/g\to0$ (\emph{strictly smaller}).
\begin{bul}
\item $f(x)-f(x_0)-K(x-x_0)=o(|x-x_0|)$ \ra\ \textbf{differentiable}. Note $(x-x_0)$, \emph{not} $Kx$: with $Kx$ the statement is false/meaningless.
\item $f(x)-f(x_0)=O(|x-x_0|)$ \ra\ only Lipschitz near $x_0$ \ra\ \textbf{not} differentiable ($|x|$ satisfies it).
\item Only $h\to0^+$ exists \ra\ only \emph{right}-differentiable \ra\ \textbf{not} differentiable.
\item "for an arbitrary polynomial $p$" \ra\ useless, $p$ must be the linear Taylor part.
\end{bul}
}
\R{piecewise function: differentiable at the seam $x_0$?}{
\begin{rcp}
\item Continuity first: $\lim_{x\to x_0^-}f=\lim_{x\to x_0^+}f=f(x_0)$. Fails \ra\ not differentiable, done.
\item Compute the one-sided \emph{difference quotients} (not just $\lim f'$) and compare. Equal \ra\ differentiable with that value.
\item For $C^1$: additionally $\lim_{x\to x_0^-}f'(x)=\lim_{x\to x_0^+}f'(x)$.
\end{rcp}
Ex.\ $f(x)=x^2\sin\frac1x$ ($x\ne0$), $f(0)=0$: $f'(0)=\lim h\sin\frac1h=0$ exists, but $f'$ is not continuous at $0$ \ra\ diff.\ but not $C^1$.
}

\subsection{Rules --- and how not to misuse them}
\F{the five rules}{
$(f\pm g)'=f'\pm g'$;\quad $(fg)'=f'g+fg'$;\quad
$\left(\frac fg\right)'=\frac{f'g-fg'}{g^2}$;\quad
\hd{chain} $(g\circ f)'(x)=g'(f(x))\cdot f'(x)$;\quad
\hd{inverse} $(f^{-1})'(y_0)=\frac1{f'(x_0)}$ with $y_0=f(x_0)$.\\
\hd{Logarithmic differentiation} for $u(x)^{v(x)}$: write $u^v=e^{v\ln u}$ \ra\ $(u^v)'=u^v\left(v'\ln u+\frac{vu'}{u}\right)$.
}
\X{using the derivative/antiderivative table (\S9) correctly}{
The table lists $F(x)$ for the \emph{bare} argument $x$. If the argument is $ax+b$:
\[\int f(ax+b)\dx=\tfrac1a F(ax+b)+C\qquad\text{(divide by the inner derivative)}\]
\warn This works \textbf{only for linear inner functions}. For anything else substitute properly.\\
\hd{Worked trap:} $\int\frac{1}{2\cos^2(\frac12x)}\dx$. Table: $\int\frac{1}{\cos^2 u}du=\tan u$.
Inner $u=\frac12x$ \ra\ $u'=\frac12$ \ra\ divide by $\frac12$, i.e.\ \emph{multiply} by 2, and keep the prefactor $\frac12$:
$\frac12\cdot 2\tan(\frac12 x)=\tan\!\left(\frac x2\right)+C$.
\hd{Always verify by differentiating your answer} --- 5 seconds, catches every chain-rule slip.
}

\subsection{Mean value theorems \de{Mittelwerts\"atze}}
\F{Fermat / Rolle / Lagrange / Cauchy}{
\begin{bul}
\item \hd{Fermat:} $x_0$ \emph{interior} local extremum and $f$ diff.\ there \ra\ $f'(x_0)=0$. (Converse false: $x^3$.)
\item \hd{Rolle:} $f$ cont.\ on $[a,b]$, diff.\ on $(a,b)$, $f(a)=f(b)$ \ra\ $\exists\xi\in(a,b):f'(\xi)=0$.
\item \hd{MVT (Lagrange):} same hypotheses \ra\ $\exists\xi\in(a,b):f'(\xi)=\frac{f(b)-f(a)}{b-a}$.
\item \hd{Cauchy MVT:} $\frac{f(b)-f(a)}{g(b)-g(a)}=\frac{f'(\xi)}{g'(\xi)}$ (this is what proves l'Hospital).
\end{bul}
}
\R{typical MVT uses}{
\begin{bul}
\item \hd{Prove an inequality} $|f(x)-f(y)|\le L|x-y|$: MVT + bound $|f'|\le L$.
\item \hd{Show $f$ is constant:} $f'\equiv0$ on an interval \ra\ MVT \ra\ constant.
\item \hd{Count zeros:} $n$ zeros of $f$ \ra\ $\ge n-1$ zeros of $f'$ (Rolle between consecutive zeros). Contrapositive: $f'$ has $\le k$ zeros \ra\ $f$ has $\le k+1$ zeros.
\item \hd{Compare two functions:} $g:=f-h$, show $g(a)=0$ and $g'>0$ \ra\ $f>h$ on $(a,\infty)$.
\end{bul}
}

\subsection{Curve discussion \de{Kurvendiskussion} / extrema}
\R{find and classify all extrema of $f$ on $I$}{
\begin{rcp}
\item \hd{Domain} + points where $f$ is not differentiable (kinks, poles, boundary of $I$) --- these are extremum candidates too.
\item \hd{Critical points:} solve $f'(x)=0$.
\item \hd{Classify:} $f''(x_0)>0$ \ra\ local \textbf{min}; $f''(x_0)<0$ \ra\ local \textbf{max}; $f''(x_0)=0$ \ra\ inconclusive, use the \emph{sign change of $f'$} ($-\!\to\!+$ min, $+\!\to\!-$ max, no change \ra\ saddle/inflection).
\item \hd{Global on $[a,b]$:} compare $f$ at all critical points \emph{and} at $a,b$. Largest value = global max (exists by Weierstrass).
\item \hd{Global on open/unbounded:} additionally take $\lim_{x\to\partial I}f$; if the limit exceeds all candidate values, the extremum is not attained.
\item \hd{Monotonicity:} sign table of $f'$. \hd{Convexity:} sign table of $f''$; sign change of $f''$ = \textbf{inflection point} \de{Wendepunkt}.
\item Asymptotes: $\lim_{x\to\pm\infty}\left(f(x)-(mx+c)\right)=0$ with $m=\lim\frac{f(x)}x$, $c=\lim(f(x)-mx)$.
\end{rcp}
}
\F{convexity (konvex) --- exam favourite}{
$f$ convex on $I$ $\iff$ $f''\ge0$ $\iff$ $f'$ increasing $\iff$ $f(\lambda x+(1-\lambda)y)\le\lambda f(x)+(1-\lambda)f(y)$ $\iff$ graph lies above every tangent.
\begin{bul}
\item Convex + $f'(x_0)=0$ at an interior point \ra\ $x_0$ is a \textbf{global minimum}. \emph{(True statement in MC questions.)}
\item A convex differentiable function cannot have several isolated critical points (the critical set is an interval).
\item Its global \emph{maximum} on $[a,b]$ is always attained at a \textbf{boundary} point.
\item Its global \emph{minimum} may well be interior --- so "min only on the boundary" is false.
\end{bul}
}

\E{full curve discussion of $f(x)=\dfrac{x}{1+x^2}$}{
\textbf{Domain} $\Rl$, odd. \textbf{$f'=\dfrac{1-x^2}{(1+x^2)^2}$} \ra\ zeros $x=\pm1$.\\
\textbf{$f''=\dfrac{2x(x^2-3)}{(1+x^2)^3}$} \ra\ $f''(1)=-\frac12<0$ \ra\ \textbf{max} $f(1)=\frac12$; $f''(-1)>0$ \ra\ \textbf{min} $f(-1)=-\frac12$.\\
\textbf{Inflection} $x=0,\pm\sqrt3$ (sign of $f''$ changes). \textbf{Monotone} incr.\ on $[-1,1]$, decr.\ outside.
\textbf{Asymptote} $y=0$ since $f\to0$ for $x\to\pm\infty$. Global max/min $=\pm\frac12$ (attained).
}

\subsection{l'Hospital \de{Bernoulli--de l'Hospital}}
\R{use it safely}{
\begin{rcp}
\item \hd{Check the form first:} only $\frac00$ or $\frac{\pm\infty}{\pm\infty}$. Writing it down without checking costs points.
\item Differentiate \emph{numerator and denominator separately} (\warn not the quotient rule).
\item Repeat if the form persists. If it never resolves or loops ($\frac{e^x}{e^x}$-style, $\frac{x}{\sqrt{1+x^2}}$), switch method.
\item Verify the new limit exists; if it does not, l'Hospital says \textbf{nothing} (e.g.\ $\frac{x+\sin x}{x}$).
\end{rcp}
\hd{Reduce other forms first:}
$0\cdot\infty$: $fg=\frac{f}{1/g}$.\quad
$\infty-\infty$: common denominator or conjugate.\quad
$1^\infty,0^0,\infty^0$: $u^v=e^{v\ln u}$, do l'Hospital in the exponent, exponentiate at the end.
}

\E{the indeterminate forms, one line each}{
\begin{bul}
\item $\lim_{x\to0}\frac{1-\cos x}{x^2}\overset{\text{Taylor}}{=}\lim\frac{x^2/2}{x^2}=\frac12$.
\item $\lim_{x\to0^+}x\ln x=\lim\frac{\ln x}{1/x}\overset{\infty/\infty}{=}\lim\frac{1/x}{-1/x^2}=\lim(-x)=0$.
\item $\lim_{x\to0^+}x^x=e^{\lim x\ln x}=e^0=1$.
\item $\lim_{x\to\infty}\left(1+\frac ax\right)^x=e^{\lim x\ln(1+a/x)}=e^{a}$.
\item $\lim_{x\to0}\left(\frac1{\sin x}-\frac1x\right)=\lim\frac{x-\sin x}{x\sin x}=\lim\frac{x^3/6}{x^2}=0$.
\item $\lim_{x\to\infty}\frac{x+\sin x}{x}=1$ --- \warn l'Hospital gives $1+\cos x$ (no limit) and is \emph{inapplicable}; just divide by $x$.
\end{bul}
}

\subsection{Taylor --- the best limit tool}
\F{Taylor formula}{
$f(x)=\sum_{k=0}^{n}\frac{f^{(k)}(x_0)}{k!}(x-x_0)^k+R_n(x)$, \hd{Lagrange remainder}
$R_n(x)=\frac{f^{(n+1)}(\xi)}{(n+1)!}(x-x_0)^{n+1}$ for some $\xi$ between $x_0$ and $x$; also $R_n=o(|x-x_0|^n)$.\\
\hd{Error bound:} $|R_n(x)|\le\frac{\max_\xi|f^{(n+1)}(\xi)|}{(n+1)!}|x-x_0|^{n+1}$ --- that is the whole "estimate the error" task.
}
\R{limits $x\to0$ via Taylor (faster and safer than l'Hospital)}{
\begin{rcp}
\item Look at the denominator's power $x^m$. You need every series up to order $m$ (one more if things cancel).
\item Insert the standard series (\S9), expand, cancel.
\item The surviving lowest-order term gives the limit; everything else is $o(x^m)$.
\end{rcp}
}
\E{$\displaystyle\lim_{x\to0}\frac{\sin x-x+\frac{x^3}6}{x^5}$}{
$\sin x=x-\frac{x^3}{6}+\frac{x^5}{120}-\frac{x^7}{5040}+\dots$\\
Numerator $=\frac{x^5}{120}+O(x^7)$ \ra\ limit $=\boxed{\tfrac1{120}}$.\\
(l'Hospital would need \textbf{five} differentiations --- don't.)
}
\R{"approximate $f(x)$ and bound the error"}{
\begin{rcp}
\item Choose $x_0$ near $x$ where you know $f$ exactly ($0$, $1$, $\frac\pi6$, \dots) and the order $n$.
\item Write $T_n(x)=\sum_{k\le n}\frac{f^{(k)}(x_0)}{k!}(x-x_0)^k$ and evaluate.
\item Error: $|R_n|\le\frac{M}{(n+1)!}|x-x_0|^{n+1}$ with $M:=\max|f^{(n+1)}|$ on the interval between $x_0$ and $x$.
\item If a target accuracy is given, solve that bound $<\eps$ for $n$.
\end{rcp}
\hd{Ex.} $e^{0.1}$, $n=3$: $T_3=1+0.1+0.005+\frac{0.001}{6}\approx1.10517$; $M=e^{0.1}<1.2$ \ra\ $|R_3|\le\frac{1.2}{24}10^{-4}=5\cdot10^{-6}$.
}
\F{Newton's method (if it appears)}{
$x_{n+1}=x_n-\frac{f(x_n)}{f'(x_n)}$; converges quadratically near a simple root if $f'\ne0$ and $f\in C^2$.
Geometrically: the zero of the tangent at $x_n$.
}
\F{Taylor series to know by heart (all around $x_0=0$)}{
$e^x=\sum\frac{x^n}{n!}$ ($R=\infty$);\quad
$\sin x=x-\frac{x^3}{3!}+\frac{x^5}{5!}\mp$;\quad
$\cos x=1-\frac{x^2}{2!}+\frac{x^4}{4!}\mp$;\\
$\ln(1+x)=x-\frac{x^2}2+\frac{x^3}3\mp$ ($|x|<1$);\quad
$\frac1{1-x}=\sum x^n$;\quad
$(1+x)^\alpha=\sum\binom{\alpha}{n}x^n$;\\
$\arctan x=x-\frac{x^3}3+\frac{x^5}5\mp$;\quad
$\sinh x=x+\frac{x^3}{3!}+\dots$;\quad $\cosh x=1+\frac{x^2}{2!}+\dots$
\hd{Tricks:} substitute ($e^{-x^2}$), differentiate/integrate termwise, multiply series --- do \emph{not} recompute derivatives.
}
